\subsection{Data}
\label{sec:data}

The empirical foundation of this study is a dataset of 424 United States-listed equities and exchange-traded funds spanning ten years of daily trading data (January~3, 2014 to December~31, 2024).
We focus the primary single-asset analysis on the SPDR S\&P~500 ETF Trust (SPY), which tracks the S\&P~500 index and is the most liquid equity ETF by trading volume.
The data include daily open, high, low, and close prices along with volume and volume-weighted average price (VWAP) metrics, sourced from standard financial data providers.
The in-sample (IS) dataset comprises $T_{\text{IS}} = 2{,}766$ trading days.
A subsequent $T_{\text{OoS}} = 219$ trading days in 2025 are reserved as a held-out out-of-sample (OoS) test set to evaluate generalization.

Cross-asset generalization is assessed on five additional tickers spanning distinct risk profiles: NVDA (NVIDIA Corporation; high-beta technology), JNJ (Johnson \& Johnson; low-beta healthcare), JPM (JPMorgan Chase \& Co.; moderate-beta financials), AAPL (Apple Inc.; large-cap consumer technology), and QQQ (Invesco QQQ Trust; Nasdaq-100 ETF).
For the univariate fidelity evaluation each asset is fitted independently using the same training/test split.
For the cross-asset extension of Section~\ref{sec:cross_asset_methods} we additionally stack the six-asset in-sample return matrix $\mathbf{R}_{\text{IS}} \in \mathbb{R}^{T_{\text{IS}} \times 6}$ to fit the cross-sectional dependence models.

To extend the framework beyond equity returns, we additionally apply it to the CBOE Volatility Index (VIX) \cite{whaley2009vix}, which measures the 30-day implied volatility of S\&P~500 options and is widely regarded as the market's ``fear gauge.''
VIX daily close prices spanning approximately 2006--2024 constitute the in-sample dataset, with 2025 data reserved for out-of-sample testing.
VIX log returns exhibit statistical properties that differ markedly from equity returns, including positive skewness (reflecting asymmetric volatility spikes), pronounced mean reversion, and extreme excess kurtosis, making it a challenging test case that probes the generality of the CHMM framework.

\subsection{Excess Growth Rate Calculation}
\label{sec:growth_rate}

We compute the annualized excess growth rate $G_{i,j}$ for ticker $i$ between trading days $j-1$ and $j$ given the equity price series $P_{i,j}$:
\begin{equation}
    G_{i,j} \equiv \frac{1}{\Delta t} \cdot \ln\!\left(\frac{P_{i,j}}{P_{i,j-1}}\right) - r_f
    \label{eq:growth_rate}
\end{equation}
where $\Delta t = 1/252$ (one trading day expressed in years) and $r_f$ is a constant risk-free rate set to zero throughout this study.
The resulting excess growth rate has units of year$^{-1}$ and is time-additive under continuous compounding.
The annualization by $1/\Delta t$ scales the daily log-returns to an annualized basis, producing values with magnitudes on the order of $\pm 2$ standard deviations for typical daily moves.

\subsection{Continuous Hidden Markov Model}
\label{sec:chmm_def}

We define the continuous hidden Markov model (CHMM) as the tuple $\mathcal{M} = (\mathcal{S}, \mathcal{O}, \mathbf{T}, \mathbf{B}, \boldsymbol{\pi})$ \cite{rabiner1986introduction}, where
\begin{itemize}
    \item $S_t \in \mathcal{S} = \{1, 2, \ldots, K\}$ is the hidden state at time $t$, representing the prevailing market regime;
    \item $\mathcal{O} \subseteq \mathbb{R}$ is the continuous observation space of excess growth rates;
    \item $\mathbf{T} \in \mathbb{R}^{K \times K}$ is the state transition probability matrix with entries $T_{ij} = \Prob(S_{t+1} = j \mid S_t = i)$, satisfying $\sum_{j=1}^{K} T_{ij} = 1$ for all $i$;
    \item $\mathbf{B} = \{(\mu_k, \sigma_k)\}_{k=1}^{K}$ specifies the Gaussian emission distributions, where $G_t \mid S_t = k \sim \mathcal{N}(\mu_k, \sigma_k^2)$;
    \item $\boldsymbol{\pi} = (\pi_1, \ldots, \pi_K)$ is the initial state distribution with $\pi_k = \Prob(S_1 = k)$.
\end{itemize}

The emission density for state $k$ is
\begin{equation}
    b_k(x) = \frac{1}{\sigma_k \sqrt{2\pi}} \exp\!\left(-\frac{(x - \mu_k)^2}{2\sigma_k^2}\right).
    \label{eq:emission}
\end{equation}

The marginal distribution of $G_t$ under the stationary regime is a $K$-component Gaussian mixture,
\begin{equation}
    f(x) = \sum_{k=1}^{K} \bar{\pi}_k \, b_k(x),
    \label{eq:mixture}
\end{equation}
where $\bar{\boldsymbol{\pi}}$ is the stationary distribution satisfying $\bar{\boldsymbol{\pi}} = \bar{\boldsymbol{\pi}} \mathbf{T}$.
This mixture structure is the mechanism by which HMMs produce heavy-tailed marginal distributions from thin-tailed Gaussian components: tail states with large $|\mu_k|$ or $\sigma_k$ contribute to the kurtosis of the mixture even though each component has zero excess kurtosis individually.

Unlike the discrete framework of \citet{alswaidan2026hybrid}, which discretizes observations into Laplace quantile bins and estimates $\mathbf{T}$ by frequency counting, we learn $\mathbf{T}$, $\{\mu_k\}_{k=1}^{K}$, $\{\sigma_k\}_{k=1}^{K}$, and $\boldsymbol{\pi}$ jointly from the continuous observations via the Baum-Welch algorithm described below.

\subsection{Baum-Welch (EM) Algorithm}
\label{sec:baum_welch}

The Baum-Welch algorithm \cite{baum1970maximization} is the standard expectation-maximization (EM) procedure for maximum likelihood estimation of HMM parameters.
Given a sequence of $T$ observations $O_{1:T} = (O_1, O_2, \ldots, O_T)$, the algorithm iterates between an E-step and an M-step until convergence.

\paragraph{E-Step: Forward-Backward in Log-Space.}
Define the forward variable
$\alpha_t(k) = \Prob(O_{1:t},\, S_t = k \mid \mathcal{M})$
and the backward variable
$\beta_t(k) = \Prob(O_{t+1:T} \mid S_t = k, \mathcal{M})$.
The recursions are
\begin{align}
    \log \alpha_1(k) &= \log \pi_k + \log b_k(O_1) \label{eq:forward_init} \\
    \log \alpha_t(j) &= \underset{i}{\text{logsumexp}}\!\left(\log \alpha_{t-1}(i) + \log T_{ij}\right) + \log b_j(O_t), \quad t = 2, \ldots, T \label{eq:forward} \\[6pt]
    \log \beta_T(k) &= 0 \label{eq:backward_init} \\
    \log \beta_t(i) &= \underset{j}{\text{logsumexp}}\!\left(\log T_{ij} + \log b_j(O_{t+1}) + \log \beta_{t+1}(j)\right), \quad t = T{-}1, \ldots, 1 \label{eq:backward}
\end{align}
where the logsumexp operation is defined as
\begin{equation}
    \underset{i}{\text{logsumexp}}(x_i) = x_{\max} + \log\!\left(\sum_{i} \exp(x_i - x_{\max})\right), \quad x_{\max} = \max_i\, x_i.
    \label{eq:logsumexp}
\end{equation}
All forward-backward computations are performed entirely in log-space to prevent the floating-point underflow that occurs when multiplying many small probabilities for long observation sequences \cite{bilmes1998gentle}.

The posterior state occupancy probability is
\begin{equation}
    \gamma_t(k) = \Prob(S_t = k \mid O_{1:T}, \mathcal{M}) = \frac{\alpha_t(k)\,\beta_t(k)}{\sum_{j=1}^{K} \alpha_t(j)\,\beta_t(j)}
    \label{eq:gamma}
\end{equation}
and the posterior transition probability is
\begin{equation}
    \xi_t(i,j) = \Prob(S_t = i, S_{t+1} = j \mid O_{1:T}, \mathcal{M}) = \frac{\alpha_t(i)\,T_{ij}\,b_j(O_{t+1})\,\beta_{t+1}(j)}{\sum_{m=1}^{K}\sum_{n=1}^{K}\alpha_t(m)\,T_{mn}\,b_n(O_{t+1})\,\beta_{t+1}(n)}.
    \label{eq:xi}
\end{equation}

\paragraph{M-Step: Parameter Re-Estimation.}
Given the posterior quantities $\gamma_t(k)$ and $\xi_t(i,j)$, the parameters are re-estimated as
\begin{align}
    \pi_k^{\text{new}} &= \gamma_1(k) \label{eq:pi_update} \\[4pt]
    \mu_k^{\text{new}} &= \frac{\displaystyle\sum_{t=1}^{T} \gamma_t(k) \cdot O_t}{\displaystyle\sum_{t=1}^{T} \gamma_t(k)} \label{eq:mu_update} \\[4pt]
    \left(\sigma_k^{\text{new}}\right)^2 &= \frac{\displaystyle\sum_{t=1}^{T} \gamma_t(k) \cdot (O_t - \mu_k^{\text{new}})^2}{\displaystyle\sum_{t=1}^{T} \gamma_t(k)} \label{eq:sigma_update} \\[4pt]
    T_{ij}^{\text{new}} &= \frac{\displaystyle\sum_{t=1}^{T-1} \xi_t(i,j)}{\displaystyle\sum_{t=1}^{T-1} \gamma_t(i)}. \label{eq:T_update}
\end{align}

\paragraph{Quantile-Based Initialization.}
The EM algorithm is well known to be sensitive to initial parameter values, with poor starting configurations leading to degenerate local optima where all states collapse to the global mean \cite{rabiner1986introduction}.
We employ a quantile-based initialization strategy: the $T$ observations are sorted in ascending order and partitioned into $K$ equal-sized chunks.
The sample mean and standard deviation of each chunk seed the corresponding state's emission parameters $(\mu_k^{(0)}, \sigma_k^{(0)})$.
The transition matrix $\mathbf{T}^{(0)}$ is initialized as $T_{ij}^{(0)} = 1/K$ for all $i, j$, and the initial distribution as $\pi_k^{(0)} = 1/K$.
This strategy ensures that the initial emission distributions span the full range of observations, with tail states initialized to capture extreme returns.

\paragraph{Convergence Criterion.}
The algorithm terminates when the absolute change in log-likelihood falls below a tolerance threshold,
\begin{equation}
    |\mathcal{L}^{(n)} - \mathcal{L}^{(n-1)}| < \text{tol},
    \label{eq:convergence}
\end{equation}
where $\mathcal{L}^{(n)} = \text{logsumexp}_{k}(\log \alpha_T^{(n)}(k))$ is the total log-likelihood at iteration $n$.
We set $\text{tol} = 10^{-4}$ and $\texttt{max\_iter} = 60$.
In practice, all models converged within 20--40 iterations (see Appendix~\ref{sec:supplemental} for convergence curves).

\paragraph{Simulation.}
Given a trained model $\mathcal{M}$, synthetic paths of length $T$ are generated by:
\begin{enumerate}
    \item Draw the initial state $S_1 \sim \text{Categorical}(\boldsymbol{\pi})$.
    \item For $t = 1, 2, \ldots, T$:
    \begin{enumerate}
        \item Emit $G_t \sim \mathcal{N}(\mu_{S_t}, \sigma_{S_t}^2)$.
        \item Transition to $S_{t+1} \sim \text{Categorical}(\mathbf{T}_{S_t, \cdot})$.
    \end{enumerate}
\end{enumerate}
No jump mechanism is applied.
The Markov transition dynamics alone govern the temporal evolution of the hidden state sequence.

\subsection{Discrete HMM Baseline (Prior Paper)}
\label{sec:discrete_hmm}

For a like-for-like comparison against the prior discrete framework of \citet{alswaidan2026hybrid}, we reproduce their discrete HMM exactly within the same evaluation pipeline.
Continuous returns are binned into $K_{\text{disc}} = 13$ quantile bins using the observed IS CDF.
Let $b(G_t) \in \{1, \dots, K_{\text{disc}}\}$ be the bin index for observation $G_t$; the transition matrix is estimated by frequency counting
\begin{equation}
    \hat{T}_{ij} \;=\; \frac{\#\{t : b(G_{t-1}) = i,\, b(G_t) = j\}}{\#\{t : b(G_{t-1}) = i\}},
    \label{eq:disc_T}
\end{equation}
and emissions are near-identity (each state emits its own bin index, with a small smoothing floor to avoid zero-probability entries).
Decoded bin indices are mapped back to continuous returns via per-bin centroids.
The ``no-jump'' (NJ) variant runs the chain with this transition matrix alone; the ``with-jumps'' (WJ) variant augments the dynamics with a Poisson jump-duration process at jump probability $\epsilon = 0.01$ and mean duration $\lambda = 3$, matching the grid-calibrated values reported by \citet{alswaidan2026hybrid}.
Because quantization onto $K_{\text{disc}} = 13$ bin centroids strips the within-bin tail mass from every simulated path, the discrete baselines provide a direct upper bound on how much of the CHMM's fidelity is attributable to continuous emissions rather than to regime structure per se.

\subsection{GARCH(1,1) Benchmark}
\label{sec:garch}

As the primary temporal benchmark, we fit a GARCH(1,1) model \cite{bollerslev1986generalized}:
\begin{align}
    G_t &= \mu + \sigma_t\, z_t, \quad z_t \overset{\text{i.i.d.}}{\sim} \mathcal{N}(0,1) \label{eq:garch_mean} \\
    \sigma_t^2 &= \omega + \alpha\, (G_{t-1} - \mu)^2 + \beta\, \sigma_{t-1}^2 \label{eq:garch_var}
\end{align}
where $\omega > 0$, $\alpha \geq 0$, $\beta \geq 0$, and $\alpha + \beta < 1$.
Parameters $(\mu, \omega, \alpha, \beta)$ are estimated by maximum likelihood via the Nelder--Mead simplex algorithm.
The GARCH model serves as a benchmark because it is the canonical parametric model for volatility dynamics: its single-decay parameter $\beta$ directly encodes exponential variance persistence, providing an upper bound on ACF reproduction for single-regime models.

\subsection{Baseline Generators}
\label{sec:baselines}

We include three additional generators as distributional reference points.

\paragraph{Bootstrap Resampling.}
Each synthetic path is constructed by sampling $T$ observations uniformly with replacement from the in-sample data.
The bootstrap preserves the exact empirical marginal distribution but destroys all temporal structure, serving as the non-parametric ceiling for distributional fidelity.

\paragraph{Gaussian i.i.d.\ Generator.}
Draws $G_t \overset{\text{i.i.d.}}{\sim} \mathcal{N}(\hat{\mu}, \hat{\sigma}^2)$, where $\hat{\mu}$ and $\hat{\sigma}$ are the sample mean and standard deviation of the in-sample data.
This generator serves as a negative control: it should fail all distributional tests due to its inability to produce heavy tails.

\paragraph{Laplace i.i.d.\ Generator.}
Draws $G_t \overset{\text{i.i.d.}}{\sim} \text{Laplace}(\hat{m}, \hat{b})$, where $\hat{m}$ is the sample median and $\hat{b}$ is the mean absolute deviation from the median.
The Laplace distribution produces moderate heavy tails (excess kurtosis $= 3$), providing an intermediate reference between the Gaussian and the heavy-tailed empirical distribution.

\subsection{Cross-Asset Dependence Models}
\label{sec:cross_asset_methods}

Let $\mathbf{R}_{\text{IS}} \in \mathbb{R}^{T \times d}$ denote the in-sample return matrix for $d$ assets, with columns indexed by $j = 1, \ldots, d$ and the market-factor column indexed by $m \in \{1, \ldots, d\}$.
Both constructions below preserve each asset's fitted CHMM marginal; they differ in how cross-asset dependence is imposed.

\paragraph{Single Index Model (SIM).}
Following \citet{sharpe1963simplified} and the multi-asset extension of \citet{alswaidan2026hybrid}, for each non-market asset $j \neq m$ we fit the affine factor equation
\begin{equation}
    G_{j,t} \;=\; \alpha_j \;+\; \beta_j\, G_{m,t} \;+\; \eta_{j,t}
    \label{eq:sim}
\end{equation}
by ordinary least squares on $\mathbf{R}_{\text{IS}}$, yielding estimates $\hat{\alpha}_j$, $\hat{\beta}_j$, per-asset residuals $\hat{\eta}_{j,1:T}$, and coefficient of determination $R^2_j$.
At simulation time, we draw $n_{\text{paths}}$ independent market-factor paths $G_{m,1:T}^{(p)}$ from the market-asset CHMM, resample idiosyncratic residuals $\hat{\eta}_{j,t}^{(p)}$ with replacement from the empirical residual vector, and compute
\begin{equation}
    \hat{G}_{j,t}^{(p)} \;=\; \hat{\alpha}_j \;+\; \hat{\beta}_j\, G_{m,t}^{(p)} \;+\; \hat{\eta}_{j,t}^{(p)}.
    \label{eq:sim_sim}
\end{equation}
The market column $G_{m,t}^{(p)}$ is itself a CHMM simulation, ensuring that the univariate fidelity of the market engine is inherited by the factor construction.
SIM is parsimonious (only $2(d-1)$ parameters plus the residual empirical distribution) but distorts each non-market marginal through the affine map in equation~\eqref{eq:sim_sim}.

\paragraph{Copulas.}
By Sklar's theorem \cite{sklar1959fonctions}, any joint cumulative distribution function $H(g_1, \dots, g_d)$ with continuous marginals $F_1, \dots, F_d$ decomposes uniquely as
\begin{equation}
    H(g_1, \dots, g_d) \;=\; C\bigl(F_1(g_1), \dots, F_d(g_d)\bigr),
    \label{eq:sklar}
\end{equation}
for some copula $C : [0,1]^d \to [0,1]$.
This decomposition underpins a principled approach to multi-asset simulation: fix each marginal $F_j$ to the per-asset CHMM and estimate only the copula $C$ from the cross-sectional co-movement.

\textit{Probability integral transform.}
We obtain pseudo-uniform observations via the nonparametric PIT $u_{j,t} = \text{rank}(g_{j,t}) / (T+1)$, which avoids the marginal misspecification risk of plugging in a parametric CDF.

\textit{Gaussian copula.}
The Gaussian copula $C_{\text{Ga}}(\mathbf{u}; \Sigma) = \Phi_d\bigl(\Phi^{-1}(u_1), \dots, \Phi^{-1}(u_d); \Sigma\bigr)$ is parameterized by the $d \times d$ correlation matrix $\Sigma$.
We estimate $\Sigma$ from the pairwise Kendall's $\tau$ matrix via the closed-form inversion $\rho_{ij} = \sin(\pi \tau_{ij} / 2)$ \cite{embrechts2002correlation, mcneil2015quantitative}, and project to the nearest positive semi-definite matrix with rescaled unit diagonal if needed.
The Gaussian copula has zero upper and lower tail dependence for all $\rho < 1$, which motivates the Student-t alternative.

\textit{Student-t copula.}
The Student-t copula $C_t(\mathbf{u}; \Sigma, \nu) = t_{d,\nu}\bigl(t_\nu^{-1}(u_1), \dots, t_\nu^{-1}(u_d); \Sigma\bigr)$ admits nonzero symmetric tail dependence $\lambda_U = \lambda_L = 2\, t_{\nu+1}\!\bigl(-\sqrt{(\nu+1)(1-\rho)/(1+\rho)}\bigr)$, controlled by $\nu$ \cite{demarta2005tcopula}.
We reuse the Kendall's-$\tau$ correlation estimate from the Gaussian copula and select $\nu$ by profile maximum likelihood over the discrete grid $\nu \in \{2, 3, 4, 5, 6, 8, 10, 15, 20, 30\}$.

\textit{Simulation via rank reordering.}
Following the rank-coupling construction of \citet{iman1982distribution}, cross-asset simulation proceeds path-by-path:
\begin{enumerate}
    \item Simulate $T$ independent paths $\tilde{g}_{j,1:T}^{(p)}$ from each asset's marginal CHMM, producing a matrix $\tilde{\mathbf{G}}^{(p)} \in \mathbb{R}^{T \times d}$.
    \item Sample $T$ copula variates $\mathbf{U}^{(p)} \in [0,1]^{T \times d}$ from $C$.
    \item For each column $j$, reorder $\tilde{g}_{j,1:T}^{(p)}$ so that its ranks match the ranks of $U_{j,1:T}^{(p)}$.
\end{enumerate}
Because reordering only permutes the simulated values, the marginal distribution of each column is preserved \emph{exactly}: the resulting $\hat{G}_{j,1:T}^{(p)}$ is a permutation of $\tilde{g}_{j,1:T}^{(p)}$, so every order statistic matches.
The copula, however, is imposed through the coupling of ranks across columns.

\subsection{Validation Metrics}
\label{sec:metrics}

We evaluate synthetic data quality using seven complementary metrics, following the framework of \cite{stenger2024thinking} and \cite{alswaidan2026hybrid}.
For each model, 1{,}000 independent paths of length $T$ are simulated (200 paths for the cross-asset extension, given the quadratic cost of the copula and SIM pipelines).
Each metric is computed per path and then aggregated.

\begin{enumerate}
    \item \textbf{KS pass rate (\%).}
    The proportion of simulated paths for which the two-sample Kolmogorov--Smirnov test \cite{kolmogorov1933, smirnov1948} fails to reject the null hypothesis of distributional equivalence with the reference data at significance level $\alpha = 0.05$.
    The KS statistic measures the maximum vertical distance between two empirical CDFs:
    \begin{equation}
        D_{n,m} = \sup_x |F_n(x) - G_m(x)|.
        \label{eq:ks}
    \end{equation}
    Higher pass rates indicate better distributional fidelity.

    \item \textbf{AD pass rate (\%).}
    The proportion of paths for which the two-sample Anderson--Darling test fails to reject at $\alpha = 0.05$.
    The AD test is more sensitive to tail discrepancies than the KS test, providing a complementary measure of distributional fit.

    \item \textbf{Excess kurtosis.}
    The mean simulated excess kurtosis across paths, compared to the observed value.
    Closer matching indicates better tail reproduction.
    Given the Gaussian emission assumption, the CHMM is expected to underestimate kurtosis relative to the true heavy-tailed distribution.

    \item \textbf{ACF-MAE.}
    The mean absolute error between the observed and mean simulated autocorrelation functions of absolute excess growth rates over $L = 252$ lags (one trading year),
    \begin{equation}
        \text{ACF-MAE} = \frac{1}{L} \sum_{\tau=1}^{L} \left|\hat{\rho}^{\text{obs}}_{|G|}(\tau) - \bar{\hat{\rho}}^{\,\text{sim}}_{|G|}(\tau)\right|,
        \label{eq:acf_mae}
    \end{equation}
    where $\bar{\hat{\rho}}^{\,\text{sim}}_{|G|}(\tau)$ is the mean simulated ACF at lag $\tau$.
    Lower values indicate better reproduction of volatility clustering.
    This is the critical metric for evaluating the Ryd\'{e}n et al.\ limitation.

    \item \textbf{Wasserstein-1 distance ($W_1$).}
    The mean absolute difference of sorted empirical quantiles, equivalent to the $L^1$ distance between the empirical CDFs \cite{glasserman2003monte},
    \begin{equation}
        W_1(F, G) = \int_0^1 |F^{-1}(u) - G^{-1}(u)|\,du,
        \label{eq:wasserstein}
    \end{equation}
    approximated by comparing sorted samples.
    Lower values indicate closer distributional match.

    \item \textbf{Hellinger distance ($H$).}
    A histogram-based distributional overlap distance in $[0,1]$,
    \begin{equation}
        H(P, Q) = \frac{1}{\sqrt{2}} \sqrt{\sum_{i=1}^{B} \left(\sqrt{p_i} - \sqrt{q_i}\right)^2},
        \label{eq:hellinger}
    \end{equation}
    where $p_i$ and $q_i$ are the bin probabilities of the observed and simulated histograms, respectively.
    $H = 0$ indicates identical distributions.

    \item \textbf{Quantile coverage (\%).}
    The fraction of 99 equally spaced quantiles (1st through 99th percentiles) of the observed distribution that fall within the $[5\text{th}, 95\text{th}]$ percentile envelope of the corresponding simulated quantiles across 1{,}000 paths.
    A coverage of 100\% indicates that the simulated distribution fully envelopes the observed distribution at all quantile levels, providing a calibration-style diagnostic of the simulation envelope.
\end{enumerate}

For the cross-asset extension we additionally track two cross-sectional metrics: the mean Frobenius norm $\|\hat{\Sigma}_{\text{sim}}^{(p)} - \hat{\Sigma}_{\text{obs}}\|_F$ and the mean absolute error of the off-diagonal Pearson correlations between simulated and observed return matrices, averaged over paths.
These quantify how faithfully each dependence model reproduces the cross-sectional correlation structure.

\subsection{Volatility Measures: VIX and VXX as Synthetic-Data Targets}
\label{sec:vix_methods}

To demonstrate that the same Baum-Welch pipeline generalizes beyond equities, we apply it to the two most widely tracked market volatility measures: the CBOE Volatility Index (VIX) and the iPath Series B S\&P 500 VIX Short-Term Futures ETN (VXX).
Both are treated as univariate return series and fitted independently.
The model specification, training, and simulation are identical to the equity pipeline of Sections~\ref{sec:chmm_def}--\ref{sec:baum_welch}; only the input series changes.
Daily log returns are computed as
\begin{equation}
    R_t^{\text{VIX}} = \ln\!\left(\frac{\text{VIX}_t}{\text{VIX}_{t-1}}\right),
    \qquad
    R_t^{\text{VXX}} = \ln\!\left(\frac{\text{VXX}_t}{\text{VXX}_{t-1}}\right).
    \label{eq:vix_returns}
\end{equation}

The same seven-metric evaluation from Section~\ref{sec:metrics} is applied path-by-path to 1{,}000 simulated return paths for each instrument.
The goal here is synthetic-data quality on volatility measures, which have positive skewness, pronounced mean reversion, and heavier-than-equity excess kurtosis; downstream risk applications that consume synthetic volatility paths (for example, VaR of volatility-sensitive portfolios or scenario stress tests) are direct consumers of this output.
Option pricing, implied-volatility-surface calibration, and stochastic-volatility benchmarking are explicitly out of scope and are not revisited in this paper.
